% !TEX root = ../main.tex
Il file \textit{plot\_utils.py} raccoglie le funzioni di visualizzazione utilizzate per rappresentare graficamente i dati sperimentali. Per garantire accessibilità, tutti i colori usati nei grafici sono scelti da una palette \textit{colorblind-safe} e definiti come costanti globali:

\begin{itemize}
    \item \texttt{C\_BAR = "\#4878CF"}: blu, usato per le barre dell'istogramma;
    \item \texttt{C\_MEAN = "\#D65F5F"}: rosso, usato per la linea della media;
    \item \texttt{C\_BAND\_A = "\#4878CF"}: blu, usato per la banda $\sigma_A$;
    \item \texttt{C\_BAND\_B = "\#EE854A"}: arancione, usato per la banda $\sigma_\text{tot}$.
\end{itemize}

Le funzioni disponibili sono:
\begin{itemize}
    \item \texttt{histogram(x, ...)}: genera un istogramma dei dati, con linea della media e banda di deviazione standard opzionali.
\end{itemize}

\section{histogram(x, bins, xlabel, ylabel, title, figsize, save\_path, show\_mean, show\_std, dpi, ax)}

\begin{description}
  \item[\textbf{Parametri}]
    \texttt{x}: array-like (lista, tupla, array NumPy) $\to$ convertito in \texttt{float64};\\
    \texttt{bins}: intero o \texttt{"auto"} (default: \texttt{"auto"}) $\to$ numero di intervalli dell'istogramma;\\
    \texttt{xlabel}: stringa (default: \texttt{"Valore1"}) $\to$ etichetta dell'asse $x$;\\
    \texttt{ylabel}: stringa o \texttt{None} (default: \texttt{None}) $\to$ etichetta dell'asse $y$;\\
    \texttt{title}: stringa (default: \texttt{"Istogramma"}) $\to$ titolo del grafico;\\
    \texttt{figsize}: tupla (default: \texttt{(8, 5)}) $\to$ dimensioni della figura in pollici;\\
    \texttt{save\_path}: stringa o \texttt{None} (default: \texttt{None}) $\to$ percorso in cui salvare la figura;\\
    \texttt{show\_mean}: \texttt{bool} (default: \texttt{True}) $\to$ mostra la linea verticale della media;\\
    \texttt{show\_std}: \texttt{bool} (default: \texttt{True}) $\to$ mostra la banda della deviazione standard;\\
    \texttt{dpi}: intero (default: \texttt{300}) $\to$ risoluzione della figura;\\
    \texttt{ax}: \texttt{matplotlib.axes.Axes} o \texttt{None} (default: \texttt{None}) $\to$ asse su cui disegnare.
  \item[\textbf{Vincoli}]
    Se \texttt{ax} $\neq$ \texttt{None}, la figura viene disegnata sull'asse passato senza crearne una nuova.
    Se \texttt{save\_path} $\neq$ \texttt{None}, la figura viene salvata nel percorso specificato con risoluzione \texttt{dpi}.
  \item[\textbf{Restituisce}]
    \texttt{int}: \texttt{0} (implementazione in corso).
  \item[\textbf{Errori}]
    Nessuno.
\end{description}

La funzione \textbf{histogram(x, ...)} genera un istogramma dei dati contenuti in \texttt{x}. Se non viene passato un asse \texttt{ax}, crea una nuova figura con le dimensioni e la risoluzione specificate.
\begin{minted}{python}
def histogram(
    x,
    bins="auto",
    xlabel="Valore1",
    ylabel=None,
    title="Istogramma",
    figsize=(8, 5),
    save_path=None,
    show_mean=True,
    show_std=True,
    dpi=300,
    ax=None,
):
    x = np.asarray(x, dtype=float)

    if ax is None:
        fig, ax = plt.subplots(figsize=figsize, dpi=dpi, constrained_layout=True)
    else:
        fig = ax.get_figure()
\end{minted}

Il primo passo converte l'input in un array \texttt{float64}. La gestione dell'asse permette di integrare la funzione sia in figure autonome sia in griglie di sottografici: se \texttt{ax = None} viene creata una nuova figura, altrimenti si usa l'asse fornito e si recupera la figura padre con \texttt{ax.get\_figure()}.